% --- Langue et polices ------------------------------------------------
\usepackage[french]{babel}
\usepackage{fontspec}

% Polices fournies avec TeX Live (aucune installation requise) :
% Source Sans (texte), Source Code (code), Lete Sans Math (mathématiques,
% sans empattements, de proportions proches de Source Sans)
\usepackage{amsmath}
\usepackage{unicode-math}
\setmainfont{SourceSansPro}[
  Extension      = .otf,
  UprightFont    = *-Regular,
  BoldFont       = *-Bold,
  ItalicFont     = *-RegularIt,
  BoldItalicFont = *-BoldIt,
]
\setmonofont{SourceCodePro}[
  Extension      = .otf,
  UprightFont    = *-Regular,
  BoldFont       = *-Bold,
  ItalicFont     = *-RegularIt,
  BoldItalicFont = *-BoldIt,
  Scale          = MatchLowercase,
]
\setmathfont{LeteSansMath.otf}

% --- Mise en page -----------------------------------------------------
\usepackage[letterpaper, margin=2cm, top=2.5cm, bottom=2.5cm, headheight=14pt]{geometry}
\setlength{\parindent}{0pt}
\setlength{\parskip}{0.5em}

% --- Paquets utilitaires ----------------------------------------------
\usepackage{xcolor}
\usepackage{graphicx}
\graphicspath{{img/}}
\usepackage{array}
\usepackage{tabularx}
\usepackage{booktabs}          % tableaux sans filets verticaux
\usepackage{enumitem}
\usepackage{multicol}
\usepackage{needspace}
\usepackage{lastpage}          % « Page x de y »
\usepackage{fancyhdr}
\usepackage[most]{tcolorbox}   % encadrés du corrigé
\usepackage{listings}          % code source

\setlist{itemsep=0.2em, topsep=0.3em}

% --- Fichiers du gabarit ----------------------------------------------
% =====================================================================
%  Le même fichier main.tex produit trois documents :
%    questionnaire       énoncés et choix de réponses
%    cahier de réponses  grille QCM et lignes de réponse, sans les énoncés
%    corrigé             le cahier de réponses, page pour page, avec les
%                        solutions
%
%  Le mode est déduit du nom du fichier compilé (voir le dossier cibles/) :
%  il doit contenir « questionnaire », « cahier_reponses » ou « corrige ».
%
%  Ce fichier fournit les commandes LaTeX suivantes :
%    \OnlyQuestionnaire{...}  seulement dans le questionnaire
%    \OnlyAnswerSheet{...}    seulement dans le cahier de réponses
%    \OnlyAnswerKey{...}      seulement dans le corrigé
%    \DocumentName            nom du document affiché dans l'en-tête et sur
%                             la couverture
% =====================================================================

\ExplSyntaxOn

\str_new:N \g_exam_mode_str
\str_gset:Nn \g_exam_mode_str { questionnaire }

\AddToHook { begindocument / before }
  {
    \str_if_in:NnTF \c_sys_jobname_str { corrige }
      { \str_gset:Nn \g_exam_mode_str { answerkey } }
      {
        \str_if_in:NnTF \c_sys_jobname_str { cahier_reponses }
          { \str_gset:Nn \g_exam_mode_str { answersheet } }
          {
            \str_if_in:NnT \c_sys_jobname_str { questionnaire }
              { \str_gset:Nn \g_exam_mode_str { questionnaire } }
          }
      }
    \typeout { Examen~:~mode~«~\g_exam_mode_str\space~» }
  }

% --- Tests utilisés par le gabarit ------------------------------------
\prg_new_conditional:Npnn \exam_if_mode:n #1 { p, T, F, TF }
  {
    \str_if_eq:VnTF \g_exam_mode_str {#1}
      { \prg_return_true: } { \prg_return_false: }
  }

% Cahier ou corrigé : même mise en page (feuille de réponses)
\prg_new_conditional:Npnn \exam_if_answer_layout: { p, T, F, TF }
  {
    \bool_lazy_or:nnTF
      { \exam_if_mode_p:n { answersheet } } { \exam_if_mode_p:n { answerkey } }
      { \prg_return_true: } { \prg_return_false: }
  }

% --- Commandes offertes dans les questions ---------------------------
%  \OnlyQuestionnaire{...}  seulement dans le questionnaire
%  \OnlyAnswerSheet{...}    seulement dans le cahier de réponses
%  \OnlyAnswerKey{...}      seulement dans le corrigé
\NewDocumentCommand \OnlyQuestionnaire { +m } { \exam_if_mode:nT { questionnaire } {#1} }
\NewDocumentCommand \OnlyAnswerSheet   { +m } { \exam_if_mode:nT { answersheet } {#1} }
\NewDocumentCommand \OnlyAnswerKey     { +m } { \exam_if_mode:nT { answerkey } {#1} }

% Nom du document affiché dans l'en-tête et sur la couverture
\NewExpandableDocumentCommand \DocumentName { }
  {
    \str_case:Vn \g_exam_mode_str
      {
        { questionnaire } { Questionnaire }
        { answersheet }   { Cahier~de~réponses }
        { answerkey }     { Corrigé }
      }
  }

\ExplSyntaxOff

% =====================================================================
%  code.tex — Coloration syntaxique (listings), pensée pour le noir et blanc
%
%  Styles disponibles : python, javascript, nodejs, html, css, java,
%                       dockerfile, yaml
%
%  Utilisation dans une question :
%
%    \begin{code}{python}              % bloc de code
%    print("Bonjour")
%    \end{code}
%
%    \begin{code}[numbers=none]{css}   % options listings entre crochets
%    ...
%    \end{code}
%
%    \codefile{java}{questions/code/Compte.java}      % depuis un fichier
%    \codefile[firstline=3, lastline=10]{java}{...}   % extrait de lignes
%
%    \codeinline{x = 3}               % code dans le texte
%
%  Règles d'impression : mots-clés en gras, commentaires en italique gris,
%  chaînes en gris foncé.
% =====================================================================

% --- Teintes de gris (noms requis : « ! » pose problème dans listings) -
\colorlet{codecomment}{black!60}
\colorlet{codestring}{black!75}
\colorlet{codenumber}{black!55}
\colorlet{codeframe}{black!45}
\colorlet{codearrow}{black!50}
\colorlet{codeattribute}{black!70}

% --- Style commun -----------------------------------------------------
\lstdefinestyle{base}{
  basicstyle        = \ttfamily\small,
  keywordstyle      = \bfseries,
  commentstyle      = \itshape\color{codecomment},
  stringstyle       = \color{codestring},
  numbers           = left,
  numberstyle       = \scriptsize\color{codenumber},
  numbersep         = 8pt,
  xleftmargin       = 2.2em,
  framexleftmargin  = 0.3em,
  frame             = single,
  framerule         = 0.4pt,
  rulecolor         = \color{codeframe},
  columns           = fullflexible,
  keepspaces        = true,
  showstringspaces  = false,
  tabsize           = 4,
  upquote           = true,
  extendedchars     = true,
  breaklines        = true,
  postbreak         = \mbox{\textcolor{codearrow}{$\hookrightarrow$}\space},
  aboveskip         = 0.6em,
  belowskip         = 0.6em,
}

% --- Définitions de langages absents de listings ---------------------
\lstdefinelanguage{JavaScript}{
  morekeywords = {async, await, break, case, catch, class, const, continue,
    default, delete, do, else, export, extends, finally, for, from,
    function, if, import, in, instanceof, let, new, of, return, static,
    super, switch, this, throw, try, typeof, var, void, while, yield,
    true, false, null, undefined},
  sensitive    = true,
  morecomment  = [l]{//},
  morecomment  = [s]{/*}{*/},
  morestring   = [b]',
  morestring   = [b]",
  morestring   = [b]`,
}

\lstdefinelanguage{CSS}{
  alsoletter   = {-@!},
  morekeywords = {color, background, background-color, background-image,
    border, border-radius, border-color, border-width, border-style,
    margin, margin-top, margin-right, margin-bottom, margin-left,
    padding, padding-top, padding-right, padding-bottom, padding-left,
    width, height, max-width, min-width, max-height, min-height,
    display, position, top, right, bottom, left, z-index, float, clear,
    font, font-family, font-size, font-weight, font-style, line-height,
    text-align, text-decoration, text-transform, letter-spacing,
    flex, flex-direction, flex-wrap, justify-content, align-items, gap,
    grid, grid-template-columns, grid-template-rows, grid-area,
    opacity, overflow, visibility, cursor, transition, transform,
    animation, box-shadow, box-sizing, content, list-style},
  morekeywords = [2]{@media, @import, @keyframes, @font-face, !important},
  sensitive    = false,
  morecomment  = [s]{/*}{*/},
  morestring   = [b]",
  morestring   = [b]',
}

\lstdefinelanguage{HTML5}[]{HTML}{
  morekeywords = {header, nav, main, section, article, aside, footer,
    figure, figcaption, video, audio, source, canvas, template,
    button, label, select, option, textarea, output, progress,
    charset, lang, placeholder, required, defer, async, type,
    for, data, alt, src, href, rel, id, class, name, value},
  sensitive    = false,
}

% Instructions en majuscules (convention) : « from » dans un RUN reste normal
\lstdefinelanguage{Dockerfile}{
  morekeywords = {FROM, AS, RUN, CMD, ENTRYPOINT, COPY, ADD, WORKDIR, ENV,
    ARG, EXPOSE, VOLUME, USER, LABEL, HEALTHCHECK, SHELL, STOPSIGNAL,
    ONBUILD},
  sensitive    = true,
  morecomment  = [l]{\#},
  morestring   = [b]",
  morestring   = [b]',
}

% Clés en gras : tout est gras jusqu'au premier « : » de la ligne,
% puis le reste de la ligne (la valeur) revient au style normal
\lstdefinelanguage{YAML}{
  morekeywords = {true, false, null, yes, no, on, off},
  sensitive    = false,
  morecomment  = [l]{\#},
  morestring   = [b]",
  morestring   = [b]',
  moredelim    = **[l][\mdseries]{:},
}

% --- Styles par langage ----------------------------------------------
\lstdefinestyle{python}{
  style        = base,
  language     = Python,
  morekeywords = {True, False, None, async, await, nonlocal, with, as,
                  assert, yield, lambda},
}

\lstdefinestyle{javascript}{
  style    = base,
  language = JavaScript,
}

% Node.js = JavaScript + quelques identifiants propres à Node
\lstdefinestyle{nodejs}{
  style        = javascript,
  morekeywords = {require, module, exports, process},
}

\lstdefinestyle{html}{
  style            = base,
  language         = HTML5,
  tagstyle         = \bfseries,
  keywordstyle     = \color{codeattribute},
  markfirstintag   = true,
  usekeywordsintag = true,
}

\lstdefinestyle{css}{
  style         = base,
  language      = CSS,
  keywordstyle  = \bfseries,
  keywordstyle  = [2]\bfseries\itshape,
}

\lstdefinestyle{java}{
  style        = base,
  language     = Java,
  morekeywords = {var, record, sealed, permits, yield},
}

\lstdefinestyle{dockerfile}{
  style    = base,
  language = Dockerfile,
}

\lstdefinestyle{yaml}{
  style        = base,
  language     = YAML,
  basicstyle   = \ttfamily\small\bfseries,
  keywordstyle = \mdseries\itshape,
  commentstyle = \mdseries\itshape\color{codecomment},
}

\lstset{style=base}

% --- Commandes pour les questions ------------------------------------
% Bloc de code :  \begin{code}[options]{langage} ... \end{code}
\lstnewenvironment{code}[2][]{\lstset{style=#2, #1}}{}

% Code depuis un fichier : \codefile[options]{langage}{chemin}
\NewDocumentCommand{\codefile}{O{} m m}{%
  \lstinputlisting[style=#2, #1]{#3}%
}

% Code dans le texte : \codeinline{...} ou \codeinline|...|
% (ne pas l'utiliser dans l'argument d'une autre commande)
\NewDocumentCommand{\codeinline}{v}{\texttt{#1}}

% =====================================================================
%  questions.tex — Parties, QCM, questions courtes, solutions et points
%
%  \exampart{Titre}
%      Ouvre une partie (A, B, C…) sur une nouvelle page. Les points de la
%      partie sont calculés automatiquement.
%
%  \begin{instructions}
%    Consignes de la partie : un ou plusieurs paragraphes.
%  \end{instructions}
%      À placer après \exampart. N'apparaît que dans le questionnaire.
%
%  \begin{mcq}[points=1]          % [breakable] : peut être coupé entre deux pages
%    Énoncé…
%    \begin{choices}[columns=2]    % columns : facultatif
%      \wrong   Réponse fausse
%      \correct Bonne réponse
%    \end{choices}
%  \end{mcq}
%
%  \begin{question}[points=3, lines=6]    % ou box=6cm : cadre vide
%    Énoncé…     % verbe d'action du niveau de Bloom visé : Expliquez…
%    \begin{solution}
%      Réponse attendue : imprimée dans le corrigé, par-dessus les
%      lignes ou dans le cadre (prévoir assez d'espace)
%    \end{solution}
%  \end{question}
%
%  Le corrigé a exactement la mise en page du cahier : mêmes pages,
%  même grille (bonnes cases noircies), mêmes espaces de réponse.
%
%  Points et totaux : il faut compiler deux fois (latexmk le fait seul).
% =====================================================================

\makeatletter
\ExplSyntaxOn

% --- Compteurs et variables -------------------------------------------
\newcounter{question}
\newcounter{exampart}
\renewcommand{\theexampart}{\Alph{exampart}}

\fp_new:N   \g__exam_total_fp           % total de l'examen (cette compilation)
\fp_new:N   \g__exam_part_fp            % total de la partie en cours
\bool_new:N \g__exam_part_open_bool
\tl_new:N   \g__exam_part_title_tl
\int_new:N  \g__exam_part_first_int     % première question de la partie
\bool_new:N \g__exam_grid_open_bool     % grille QCM ouverte (cahier)
\int_new:N  \g__exam_nchoices_int       % nombre de choix du QCM courant
\box_new:N  \l__exam_discard_box        % boîte jetée (contenu masqué)
\box_new:N  \l__exam_mcq_box            % QCM gardé sur une seule page
\box_new:N  \g__exam_solution_box       % solution de la question courante
\seq_new:N  \g__exam_correct_seq        % numéros des bons choix du QCM
\bool_new:N \l__exam_in_question_bool
\bool_new:N \l__exam_show_solution_bool
\bool_new:N \l__exam_breakable_bool
\fp_new:N   \l__exam_points_fp
\int_new:N  \l__exam_lines_int
\dim_new:N  \l__exam_boxheight_dim
\int_new:N  \l__exam_columns_int
\bool_new:N \l__exam_multicol_bool
\tl_new:N   \l__exam_rows_tl
\skip_const:Nn \c__exam_question_skip { 2em }  % espace avant une question

\keys_define:nn { exam / question }
  {
    points .fp_set:N  = \l__exam_points_fp,
    points .initial:n = 1,
    lines     .int_set:N  = \l__exam_lines_int,
    lines     .initial:n  = 5,
    box       .dim_set:N  = \l__exam_boxheight_dim,
    box       .initial:n  = 0pt,
    breakable .bool_set:N = \l__exam_breakable_bool,
    breakable .default:n  = true,
    breakable .initial:n  = false,
  }

\keys_define:nn { exam / choices }
  {
    columns .int_set:N = \l__exam_columns_int,
    columns .initial:n = 1,
  }

% --- Valeurs lues dans le .aux (compilation précédente) ---------------
\seq_new:N  \g_exam_aux_rows_seq          % {lettre}{titre}{de}{à}{points}
\prop_new:N \g_exam_aux_part_points_prop  % lettre -> points
\tl_new:N   \g_exam_aux_total_tl

\NewDocumentCommand \ExamPartAux { m m m m m }
  {
    \seq_gput_right:Nn \g_exam_aux_rows_seq { {#1} {#2} {#3} {#4} {#5} }
    \prop_gput:Nnn \g_exam_aux_part_points_prop {#1} {#5}
  }
\NewDocumentCommand \ExamTotalAux { m }
  { \tl_gset:Nn \g_exam_aux_total_tl {#1} }

% Le .aux est relu à la fin du document : on ignore cette relecture.
\AddToHook { enddocument / afterlastpage }
  {
    \RenewDocumentCommand \ExamPartAux { m m m m m } { }
    \RenewDocumentCommand \ExamTotalAux { m } { }
  }

% --- Affichage des nombres et des points ------------------------------
% 0.5 -> 0,5 ; « 1 point », « 1,5 point », « 2 points »
\cs_new_protected:Npn \exam_print_num:n #1
  {
    \tl_set:Ne \l_tmpa_tl { \fp_to_decimal:n {#1} }
    \tl_replace_all:Nnn \l_tmpa_tl { . } { , }
    \tl_use:N \l_tmpa_tl
  }
\cs_new_protected:Npn \exam_print_points:n #1
  {
    \exam_print_num:n {#1} \nobreakspace
    \fp_compare:nNnTF {#1} < { 2 } { point } { points }
  }

\cs_new_protected:Npn \__exam_add_points:
  {
    \fp_gadd:Nn \g__exam_total_fp { \l__exam_points_fp }
    \fp_gadd:Nn \g__exam_part_fp  { \l__exam_points_fp }
  }

% --- Parties ----------------------------------------------------------
\cs_new_protected:Npn \__exam_close_part:
  {
    \bool_if:NT \g__exam_part_open_bool
      {
        \iow_now:Ne \@auxout
          {
            \token_to_str:N \ExamPartAux
              { \int_to_Alph:n { \value{exampart} } }
              { \exp_not:V \g__exam_part_title_tl }
              { \int_use:N \g__exam_part_first_int }
              { \int_use:N \c@question }
              { \fp_to_decimal:N \g__exam_part_fp }
          }
        \bool_gset_false:N \g__exam_part_open_bool
      }
  }

\cs_new_protected:Npn \__exam_part_points:
  {
    \prop_get:NeNTF \g_exam_aux_part_points_prop
      { \int_to_Alph:n { \value{exampart} } } \l_tmpb_tl
      { \exam_print_points:n { \l_tmpb_tl } }
      { ??~points }
  }
\cs_generate_variant:Nn \prop_get:NnNTF { Ne }

\NewDocumentCommand \exampart { m }
  {
    \__exam_close_grid:
    \__exam_close_part:
    \stepcounter{exampart}
    \tl_gset:Nn \g__exam_part_title_tl {#1}
    \fp_gzero:N \g__exam_part_fp
    \int_gset:Nn \g__exam_part_first_int { \value{question} + 1 }
    \bool_gset_true:N \g__exam_part_open_bool
    \clearpage   % chaque partie commence sur une nouvelle page
    \noindent
    { \Large \bfseries Partie \nobreakspace \theexampart \nobreakspace — \nobreakspace #1 }
    \hfill { \normalsize \__exam_part_points: }
    \par \nointerlineskip \vspace { 0.3em }
    \noindent \rule { \linewidth } { 1pt }
    \par
  }

% Consignes d'une partie : imprimées seulement dans le questionnaire
\NewDocumentEnvironment { instructions } { }
  { \par \exam_if_answer_layout:T { \vbox_set:Nw \l__exam_discard_box } }
  { \par \exam_if_answer_layout:T { \vbox_set_end: } }

% --- En-tête d'une question -------------------------------------------
\cs_new_protected:Npn \__exam_question_header:
  {
    \par \addvspace { \c__exam_question_skip }
    \Needspace { 5 \baselineskip }
    \__exam_question_title:
  }
\cs_new_protected:Npn \__exam_question_title:
  {
    \noindent \textbf { Question~\thequestion }
    \hfill { \small \itshape ( \exam_print_points:n { \l__exam_points_fp } ) }
    \par \nopagebreak
  }

% --- Lignes de réponse (cahier et corrigé) ----------------------------
% Dans le corrigé, les lignes sont invisibles (même espace réservé) et la
% solution est posée en haut de cet espace, sans occuper de place :
% la mise en page reste identique à celle du cahier.
\cs_new_protected:Npn \__exam_rule_line:
  {
    \exam_if_mode:nTF { answerkey }
      { \rule { 0pt } { 0.4pt } }
      { \textcolor { black!45 } { \rule { \linewidth } { 0.4pt } } }
  }

\cs_new_protected:Npn \__exam_solution_overlay:
  {
    \box_if_empty:NF \g__exam_solution_box
      {
        % haut de la solution à 0,6 cm au-dessus de la première ligne
        \dim_set:Nn \l_tmpa_dim { 0.6cm - \box_ht:N \g__exam_solution_box }
        \hbox_set:Nn \l_tmpa_box
          { \box_move_up:nn { \l_tmpa_dim } { \box_use_drop:N \g__exam_solution_box } }
        \box_set_ht:Nn \l_tmpa_box { 0pt }
        \box_set_dp:Nn \l_tmpa_box { 0pt }
        \hbox_overlap_right:n
          {
            \skip_horizontal:n { 0.5cm }
            \box_use_drop:N \l_tmpa_box
          }
      }
  }

\cs_new_protected:Npn \__exam_lines:n #1
  {
    \par
    \int_compare:nNnT {#1} > { 0 }
      {
        \group_begin:
          \skip_set:Nn \parskip { 0pt }
          \skip_set:Nn \baselineskip { 0.85cm }
          \noindent \__exam_solution_overlay: \__exam_rule_line: \par
          \prg_replicate:nn { #1 - 1 }
            { \noindent \__exam_rule_line: \par }
        \group_end:
      }
  }

% --- Cadre de réponse (cahier : vide ; corrigé : contient la solution) -
% Hauteur fixe : la mise en page est identique dans les deux documents.
\tcbset{
  exam/answerbox/.style = {
    sharp~corners, colback = white, colframe = black!60,
    boxrule = 0.6pt, before~skip = 0.5em, after~skip = 0.5em,
    left = 3mm, right = 3mm, top = 2mm, bottom = 2mm,
  }
}
\cs_new_protected:Npn \__exam_answer_box:n #1
  {
    \par
    \begin{tcolorbox} [ exam/answerbox, height = #1 ]
      \box_use_drop:N \g__exam_solution_box
    \end{tcolorbox}
  }

% --- Question courte --------------------------------------------------
%  lines=n  : n lignes de réponse dans le cahier (5 par défaut)
%  box=6cm  : cadre vide de cette hauteur, à la place des lignes
\NewDocumentEnvironment { question } { O{} }
  {
    \keys_set:nn { exam / question } {#1}
    \__exam_close_grid:
    \refstepcounter { question }
    \__exam_add_points:
    \__exam_question_header:
    \box_gclear:N \g__exam_solution_box
    \bool_set_true:N \l__exam_in_question_bool
    \exam_if_answer_layout:T { \vbox_set:Nw \l__exam_discard_box }
  }
  {
    \exam_if_answer_layout:T
      {
        \vbox_set_end:
        \dim_compare:nNnTF { \l__exam_boxheight_dim } > { 0pt }
          { \__exam_answer_box:n { \l__exam_boxheight_dim } }
          { \__exam_lines:n { \l__exam_lines_int } }
      }
    \par
  }

% --- Solution ---------------------------------------------------------
% Corrigé : mise de côté, puis placée dans l'espace de réponse de la
% question (voir \__exam_lines:n et \__exam_answer_box:n).
% Ailleurs, et dans un QCM, la solution n'est pas imprimée.
\NewDocumentEnvironment { solution } { }
  {
    \bool_set:Nn \l__exam_show_solution_bool
      { \exam_if_mode_p:n { answerkey } && \l__exam_in_question_bool }
    \bool_if:NTF \l__exam_show_solution_bool
      {
        \vbox_gset:Nw \g__exam_solution_box
        \dim_set:Nn \hsize { \linewidth - 1cm }
        \dim_set:Nn \linewidth { \hsize }
      }
      { \vbox_set:Nw \l__exam_discard_box }
  }
  {
    \bool_if:NTF \l__exam_show_solution_bool
      { \vbox_gset_end: }
      { \vbox_set_end: }
  }

% --- QCM --------------------------------------------------------------
\NewDocumentEnvironment { mcq } { O{} }
  {
    \keys_set:nn { exam / question } {#1}
    \refstepcounter { question }
    \__exam_add_points:
    \int_gzero:N \g__exam_nchoices_int
    \seq_gclear:N \g__exam_correct_seq
    \exam_if_answer_layout:TF
      { \vbox_set:Nw \l__exam_discard_box }
      {
        \bool_if:NTF \l__exam_breakable_bool
          { \__exam_question_header: }
          {
            % Par défaut, un QCM n'est pas coupé entre deux pages
            \par \addvspace { \c__exam_question_skip }
            \vbox_set:Nw \l__exam_mcq_box
            \__exam_question_title:
          }
      }
  }
  {
    \exam_if_answer_layout:TF
      {
        \vbox_set_end:
        \__exam_grid_row:
      }
      {
        \par
        \bool_if:NF \l__exam_breakable_bool
          {
            \vbox_set_end:
            \box_use_drop:N \l__exam_mcq_box
            \par
          }
      }
  }

\NewDocumentEnvironment { choices } { O{} }
  {
    \keys_set:nn { exam / choices } {#1}
    \int_gzero:N \g__exam_nchoices_int
    \setlength { \lineskip } { 0.6em }   % écart minimal entre choix très hauts
    \bool_set:Nn \l__exam_multicol_bool
      {
        \int_compare_p:nNn { \l__exam_columns_int } > { 1 }
        && ! \exam_if_answer_layout_p:
      }
    \bool_if:NT \l__exam_multicol_bool
      {
        % Première ligne de chaque colonne sur la même ligne de base, même
        % si un choix est plus haut (ex. une fraction) ; \multicolsep garde
        % un espace au-dessus d'un choix très haut (ex. une somme)
        \setlength { \multicolsep } { 0.9em }
        \setlength { \topskip } { 1.8em }
        \begin{multicols} { \l__exam_columns_int }
      }
    \begin{itemize} [ leftmargin = 2.6em, labelsep = 0.6em,
                      itemsep = 0.3em, topsep = 0.4em ]
  }
  {
    \end{itemize}
    \bool_if:NT \l__exam_multicol_bool { \end{multicols} }
  }

% Choix « A) », « B) »… ; \correct retient le numéro pour le corrigé.
\cs_new_protected:Npn \__exam_choice:
  {
    \int_gincr:N \g__exam_nchoices_int
    \use:e
      {
        \exp_not:N \item
          [ \exp_not:N \textbf { \int_to_Alph:n { \g__exam_nchoices_int } ) } ]
      }
  }
\NewDocumentCommand \wrong   { } { \__exam_choice: }
\NewDocumentCommand \correct { }
  {
    \__exam_choice:
    \seq_gput_right:Ne \g__exam_correct_seq
      { \int_use:N \g__exam_nchoices_int }
  }

% --- Grille QCM (cahier et corrigé) -----------------------------------
% Chaque QCM ajoute une ligne ; l'en-tête est ajouté au début d'une série.
% Case vide, ou noircie (même taille) pour une bonne réponse du corrigé.
\cs_new_protected:Npn \__exam_checkbox:n #1
  {
    \group_begin:
      \dim_set:Nn \fboxsep { 0pt }
      \dim_set:Nn \fboxrule { 0.6pt }
      \bool_set_false:N \l_tmpa_bool
      \exam_if_mode:nT { answerkey }
        {
          \seq_if_in:NnT \g__exam_correct_seq {#1}
            { \bool_set_true:N \l_tmpa_bool }
        }
      \raisebox { -0.15em }
        {
          \fbox
            {
              \bool_if:NTF \l_tmpa_bool
                { \rule { 0.95em } { 0.95em } }
                { \rule { 0pt } { 0.95em } \hspace* { 0.95em } }
            }
        }
    \group_end:
  }

\cs_new_protected:Npn \__exam_grid_line:nn #1#2
  {
    \penalty 0
    \nointerlineskip
    \hbox_to_wd:nn { \linewidth }
      {
        \hfil
        \begin{tabular} { | wc{6em} | wl{24em} | }
          #1 \rule [ -0.65em ] { 0pt } { 1.9em } & #2 \\ \hline
        \end{tabular}
        \hfil
      }
  }

\cs_new_protected:Npn \__exam_grid_row:
  {
    \par
    \bool_if:NF \g__exam_grid_open_bool
      {
        \bool_gset_true:N \g__exam_grid_open_bool
        \addvspace { 0.8em }
        \__exam_grid_line:nn
          { \hline \textbf { Question } }
          { \textbf { Réponse } ~ { \small ( cochez~une~seule~case ) } }
      }
    \__exam_grid_line:nn
      { \thequestion }
      {
        \int_step_inline:nn { \g__exam_nchoices_int }
          { \int_to_Alph:n {##1} \enspace \__exam_checkbox:n {##1} \hspace { 1.8em } }
      }
  }

\cs_new_protected:Npn \__exam_close_grid:
  {
    \bool_if:NT \g__exam_grid_open_bool
      {
        \bool_gset_false:N \g__exam_grid_open_bool
        \par \addvspace { 1em }
      }
  }

% --- Tableau des points (page de couverture) --------------------------
\cs_new_protected:Npn \__exam_table_row:nnnnn #1#2#3#4#5
  {
    Partie~#1~—~#2 &
    \int_compare:nNnTF {#3} > {#4} { — }
      { \int_compare:nNnTF {#3} = {#4} {#3} { #3~à~#4 } } &
    \exam_print_num:n {#5}
    \exam_if_mode:nT { answersheet } { & \hfill / \exam_print_num:n {#5} }
    \\
  }

\NewDocumentCommand \PointsTable { }
  {
    \tl_if_empty:NTF \g_exam_aux_total_tl
      { \textit { (Compilez~une~seconde~fois~pour~obtenir~le~barème.) } }
      {
        \tl_clear:N \l__exam_rows_tl
        \seq_map_inline:Nn \g_exam_aux_rows_seq
          { \tl_put_right:Nn \l__exam_rows_tl { \__exam_table_row:nnnnn ##1 } }
        \group_begin:
        \renewcommand { \arraystretch } { 1.3 }
        \use:e
          {
            \exp_not:N \begin{tabularx} { \exp_not:N \linewidth }
              { @{} >{ \exp_not:N \raggedright \exp_not:N \arraybackslash } X
                wc{6em} wc{5em}
                \exam_if_mode:nT { answersheet } { wc{6em} } @{} }
          }
          \textbf { Partie } & \textbf { Questions } & \textbf { Points }
            \exam_if_mode:nT { answersheet } { & \textbf { Résultat } } \\ \midrule
          \l__exam_rows_tl
          \midrule
          \multicolumn { 2 } { @{} r } { \textbf { Total } } &
          \textbf { \exam_print_num:n { \g_exam_aux_total_tl } }
            \exam_if_mode:nT { answersheet }
              { & \hfill \textbf { / \exam_print_num:n { \g_exam_aux_total_tl } } }
          \\
        \end{tabularx}
        \group_end:
      }
  }

% --- Fin du document --------------------------------------------------
\AddToHook { enddocument }
  {
    \__exam_close_grid:
    \__exam_close_part:
    \iow_now:Ne \@auxout
      { \token_to_str:N \ExamTotalAux { \fp_to_decimal:N \g__exam_total_fp } }
  }

\ExplSyntaxOff
\makeatother

% =====================================================================
%  couverture.tex — Paramètres de l'examen, page d'ouverture, en-têtes
%
%  Dans main.tex :
%    \ExamSetup{ title = ..., course = ..., code = ..., ... }
%    \Instructions{ \item ... \item ... }
%    \CoverPage
%
%  Contenu de la couverture selon le mode :
%    questionnaire : infos du cours, consignes, barème
%    cahier        : infos du cours, identification, consignes du cahier,
%                    barème avec colonne « Résultat »
%    corrige       : comme le questionnaire, avec la mention « Corrigé »
% =====================================================================

\ExplSyntaxOn

% --- Paramètres -------------------------------------------------------
\keys_define:nn { exam / setup }
  {
    title       .tl_gset:N = \g_exam_title_tl,
    course      .tl_gset:N = \g_exam_course_tl,
    code        .tl_gset:N = \g_exam_code_tl,
    groups      .clist_gset:N = \g_exam_groups_clist,
    teacher     .clist_gset:N = \g_exam_teacher_clist,
    teacher-feminine .bool_gset:N = \g_exam_teacher_feminine_bool,
    session     .tl_gset:N = \g_exam_session_tl,
    date        .tl_gset:N = \g_exam_date_tl,
    duration    .tl_gset:N = \g_exam_duration_tl,
    weight      .tl_gset:N = \g_exam_weight_tl,
    logo        .tl_gset:N = \g_exam_logo_tl,
    logo        .initial:n = Logo_CM_couleur_Web_grand.jpg,
    title       .initial:n = Examen,
  }
\NewDocumentCommand \ExamSetup { m }
  { \keys_set:nn { exam / setup } {#1} }

\tl_new:N \g_exam_instructions_tl
\tl_new:N \g_exam_answer_sheet_instructions_tl
\tl_gset:Nn \g_exam_answer_sheet_instructions_tl
  {
    \item Inscrivez~votre~nom~et~votre~prénom~ci-dessus.
    \item Seules~les~réponses~inscrites~dans~ce~cahier~seront~corrigées.
  }
\NewDocumentCommand \Instructions            { +m } { \tl_gset:Nn \g_exam_instructions_tl {#1} }
\NewDocumentCommand \AnswerSheetInstructions { +m } { \tl_gset:Nn \g_exam_answer_sheet_instructions_tl {#1} }

% --- En-têtes et pieds de page ----------------------------------------
% En-tête : Automne 2026 — 420-XXX-MA — Titre de l'examen (valeurs vides omises)
\seq_new:N \l__exam_header_seq
\cs_new_protected:Npn \__exam_header_id:
  {
    \seq_clear:N \l__exam_header_seq
    \tl_if_empty:NF \g_exam_session_tl { \seq_put_right:NV \l__exam_header_seq \g_exam_session_tl }
    \tl_if_empty:NF \g_exam_code_tl    { \seq_put_right:NV \l__exam_header_seq \g_exam_code_tl }
    \tl_if_empty:NF \g_exam_title_tl   { \seq_put_right:NV \l__exam_header_seq \g_exam_title_tl }
    \seq_use:Nn \l__exam_header_seq { ~—~ }
  }

\pagestyle{fancy}
\fancyhf{}
\fancyhead[L]{ \small \__exam_header_id: }
\fancyhead[R]{ \small \DocumentName }
\fancyfoot[C]{ \small Page~\thepage{}~de~\pageref{LastPage} }
\renewcommand{\headrulewidth}{0.4pt}

\fancypagestyle{cover}
  {
    \fancyhf{}
    \fancyfoot[C]{ \small Page~\thepage{}~de~\pageref{LastPage} }
    \renewcommand{\headrulewidth}{0pt}
  }

% --- Éléments de la couverture ----------------------------------------
% Titre de section de la couverture
\cs_new_protected:Npn \__exam_cover_section:n #1
  {
    \par \addvspace { 1.6em }
    \noindent { \large \bfseries #1 } \par \nopagebreak
  }

% Case du tableau d'informations (omise si la valeur est vide)
\cs_new_protected:Npn \__exam_info_row:nN #1#2
  {
    \tl_if_empty:NF #2
      { \seq_put_right:Nn \l__exam_info_seq { \textbf{#1} & #2 } }
  }
\seq_new:N \l__exam_info_seq
\tl_new:N  \l__exam_info_tl

% Case pour une liste : étiquette au pluriel (#2) s'il y a plusieurs
% éléments, sinon au singulier (#1) ; éléments joints par « , » et « et »
\cs_new_protected:Npn \__exam_info_list_row:nnN #1#2#3
  {
    \clist_if_empty:NF #3
      {
        \tl_set:Ne \l_tmpa_tl { \clist_use:Nnnn #3 { ~et~ } { ,~ } { ~et~ } }
        \int_compare:nNnTF { \clist_count:N #3 } > { 1 }
          { \tl_set:Nn \l_tmpb_tl {#2} }
          { \tl_set:Nn \l_tmpb_tl {#1} }
        \seq_put_right:Ne \l__exam_info_seq
          { \exp_not:N \textbf { \exp_not:V \l_tmpb_tl } & \exp_not:V \l_tmpa_tl }
      }
  }

\cs_new_protected:Npn \__exam_cover_info:
  {
    \seq_clear:N \l__exam_info_seq
    \bool_if:NTF \g_exam_teacher_feminine_bool
      { \__exam_info_list_row:nnN { Enseignante } { Enseignantes } \g_exam_teacher_clist }
      { \__exam_info_list_row:nnN { Enseignant }  { Enseignants }  \g_exam_teacher_clist }
    \__exam_info_list_row:nnN { Groupe } { Groupes } \g_exam_groups_clist
    \__exam_info_row:nN { Session }       \g_exam_session_tl
    \__exam_info_row:nN { Date }          \g_exam_date_tl
    \__exam_info_row:nN { Durée }         \g_exam_duration_tl
    \__exam_info_row:nN { Pondération }   \g_exam_weight_tl
    % Deux informations par ligne
    \tl_clear:N \l__exam_info_tl
    \int_zero:N \l_tmpa_int
    \seq_map_inline:Nn \l__exam_info_seq
      {
        \int_incr:N \l_tmpa_int
        \tl_put_right:Nn \l__exam_info_tl {##1}
        \int_if_odd:nTF { \l_tmpa_int }
          { \tl_put_right:Nn \l__exam_info_tl { & } }
          { \tl_put_right:Nn \l__exam_info_tl { \\ } }
      }
    \int_if_odd:nT { \l_tmpa_int }
      { \tl_put_right:Nn \l__exam_info_tl { & \\ } }
    \group_begin:
      \renewcommand { \arraystretch } { 1.3 }
      \begin{tabularx} { \linewidth } { @{} l X l X @{} }
        \l__exam_info_tl
      \end{tabularx}
    \group_end:
  }

% Champs à remplir : étiquette suivie d'une ligne
\cs_new_protected:Npn \__exam_id_row:nn #1#2
  {
    #1 & \rule { 0pt } { 2.2em } \hrulefill &
    #2 & \hrulefill \\
  }

\cs_new_protected:Npn \__exam_cover_identification:
  {
    \noindent
    \begin{tabularx} { \linewidth }
      { @{} l @{\hspace{0.5em}} X @{\hspace{2em}} l @{\hspace{0.5em}} X @{} }
      \__exam_id_row:nn { Nom } { Prénom }
    \end{tabularx}
  }

\cs_new_protected:Npn \__exam_cover_banner:
  {
    \exam_if_mode:nTF { answerkey }
      {
        \colorbox { black }
          {
            \textcolor { white }
              { \Large \bfseries \ CORRIGÉ \ }
          }
      }
      {
        { \large \bfseries \addfontfeature { LetterSpace = 12 }
          \text_uppercase:n { \DocumentName } }
      }
  }

% --- Page de couverture -----------------------------------------------
\NewDocumentCommand \CoverPage { }
  {
    \thispagestyle { cover }
    \noindent
    \includegraphics [ width = 6.5cm ] { \g_exam_logo_tl }
    \par \vspace { 0.8cm }

    \begin{center}
      { \Huge \bfseries \g_exam_title_tl \par }
      \vspace { 0.6em }
      { \LARGE \g_exam_course_tl \par }
      \tl_if_empty:NF \g_exam_code_tl
        { \vspace { 0.2em } { \large \g_exam_code_tl \par } }
      \vspace { 0.8em }
      \__exam_cover_banner:
    \end{center}

    \__exam_cover_info:

    \exam_if_mode:nTF { answersheet }
      {
        \__exam_cover_section:n { Identification }
        \__exam_cover_identification:
        \__exam_cover_section:n { Consignes }
        \begin{itemize} [ itemsep = 0pt, parsep = 0pt, topsep = 0pt ]
          \g_exam_answer_sheet_instructions_tl
        \end{itemize}
      }
      {
        \tl_if_empty:NF \g_exam_instructions_tl
          {
            \__exam_cover_section:n { Consignes }
            \begin{itemize} [ itemsep = 0pt, parsep = 0pt, topsep = 0pt ]
              \g_exam_instructions_tl
            \end{itemize}
          }
      }

    \__exam_cover_section:n { Barème }
    \vspace { 0.3em }
    \noindent \PointsTable
    \clearpage
  }

\ExplSyntaxOff

